Yes. Let $h>0$ be the distance between the parallel lines, and let the common angle be $\theta$. The perpendicular component of each segment joining the lines has length $h$, so $AC\sin\theta=h$ and $BD\sin\theta=h$. As $0<\theta<180^\circ$, its sine is positive, and division gives $\boxed{AC=BD}$. This reasoning also covers obtuse equal angles and crossing segments, since the same positive perpendicular distance is used.
\begin{asy}
import graph;import geometry;usepackage("amsmath");size(7cm);
size(12cm);for(int k=0;k<3;++k){pair v=(10*k,0),A=v,B=v+(4,0),C=v+(k==0?1:(k==1?6:-1),2),D=v+(k==0?3:(k==1?-2:5),2);draw(v+(-2.5,0)--v+(6.5,0));draw(v+(-2.5,2)--v+(6.5,2));draw(A--C^^B--D);label("$A$",A,(0,-1));label("$B$",B,(0,-1));label("$C$",C,(0,1));label("$D$",D,(0,1));}
\end{asy}

